%% ma: L12-B12-0014
%% lop: 12
%% chuong: 4
%% bai: 12
%% dang: 12.12.2
%% loai: TLN
%% muc: VD
%% dap_an: "5200"
%% nguon: "(NBV)-ĐỀ SỐ 11-MỖI NGÀY 1 ĐỀ THI 2025.pdf (D11, MỖI NGÀY 1 ĐỀ - Đề số 11), Phần III Câu 3"
%% kien_thuc: "Bài 12 (tích phân, ý nghĩa vật lí: quãng đường và vận tốc)"
%% trang_thai: chinh-thuc
%% ghi_chu: "Đáp án gốc 5200 đúng. Đề gốc có ảnh minh họa xe đua (ảnh trang trí, không chứa dữ kiện toán) nên bỏ hình. Đơn vị đáp số là mét."
\begin{ex}
Một chiếc xe đua theo công thức I bắt đầu chuyển động tăng tốc với gia tốc không đổi; khi vận tốc đạt $80$ m/s thì xe chuyển động với vận tốc không đổi trong thời gian $56$ s, sau đó giảm tốc với gia tốc không đổi đến khi xe dừng lại. Biết rằng thời gian chuyển động của xe là $74$ s. Tính quãng đường đi được của xe (đơn vị: mét).
\shortans{5200}
\loigiai{Gọi $t_1$, $t_3$ là thời gian tăng tốc và giảm tốc, $t_1+t_3=74-56=18$ (s). Vận tốc $v_1(t)=\dfrac{80}{t_1}t$ khi tăng tốc, $v_2=80$ khi chuyển động đều, $v_3(t)=80-\dfrac{80}{t_3}t$ khi giảm tốc (tính từ lúc bắt đầu giảm tốc).

Quãng đường: $S_1=\displaystyle\int_0^{t_1}\dfrac{80}{t_1}t\,\mathrm dt=40t_1$, $S_2=80\cdot56=4480$, $S_3=\displaystyle\int_0^{t_3}\left(80-\dfrac{80}{t_3}t\right)\mathrm dt=40t_3$.

$S=40(t_1+t_3)+4480=40\cdot18+4480=5200$ (m).}
\end{ex}
